\documentclass[10pt,a4paper]{amsart}
\usepackage[margin=19mm]{geometry}
\usepackage{amsmath,amssymb,mathtools}
\usepackage[scr=rsfs,cal=euler]{mathalfa}
\usepackage{xfrac}
\usepackage[unicode,hidelinks,bookmarks=false]{hyperref}
\newcommand{\ie}{i.e.\ }
\newcommand{\cf}{cf.\ }
\newcommand{\resp}{respectively\ }
\newcommand{\Subsection}{\subsection}
\providecommand{\nobreakdash}{\nobreak\hbox{-}}
\setlength{\emergencystretch}{2em}
\multlinegap0pt
\usepackage{amscd}
\usepackage{dsfont}
\usepackage{graphicx}
\numberwithin{equation}{section}
\DeclareMathOperator{\lcm}{lcm}
\DeclareMathOperator{\Arg}{Arg}
\DeclareMathOperator{\Hom}{Hom}
\DeclareMathOperator{\SM}{SM}
\DeclareMathOperator{\Herm}{Herm}
\DeclareMathOperator{\Tr}{Tr}
\DeclareMathOperator{\Pic}{Pic}
\DeclareMathOperator{\Proj}{Proj}
\DeclareMathOperator{\Vol}{Vol}
\DeclareMathOperator{\Span}{Span}
\DeclareMathOperator{\rank}{rank}
\DeclareMathOperator{\diff}{d}
\DeclareMathOperator{\dbar}{\bar\partial}
\DeclareMathOperator{\ddbar}{\partial\bar\partial}
\DeclarePairedDelimiter{\abs}{\lvert}{\rvert}
\DeclarePairedDelimiter{\Abs}{\lVert}{\rVert}
\newenvironment{equation'}{
\begin{equation}
\tag{\theequation\ensuremath{'}}}{
\end{equation}
}
\newenvironment{equation_}[1]{
\begin{equation}
\tag{\ensuremath{\theequation_{#1}}}}{
\end{equation}
}
\newenvironment{eqlist}{
\begingroup \leftskip\parindent \rightskip\parindent \noindent
\begin{minipage}{\linewidth-2\parindent}} {
\end{minipage}
\endgroup}
\newcommand*\rS{\mathrm{S}}
\newcommand*\Id{\mathrm{Id}}
\newcommand*\Sing{\mathrm{Sing}}
\newcommand*\orb{\mathrm{orb}}
\newcommand*\FS{\mathrm{FS}}
\newcommand*\loc{\mathrm{loc}}
\newcommand{\rma}{\mathrm{a}}
\newcommand{\rmb}{\mathrm{b}}
\newcommand*\cI{\mathcal{I}}
\newcommand*\cJ{\mathcal{J}}
\newcommand*\cO{\mathcal{O}}
\newcommand*\bC{\mathbb{C}}
\newcommand*\bD{\mathbb{D}}
\newcommand*\bE{\mathbb{E}}
\newcommand*\bG{\mathbb{G}}
\newcommand*\bN{\mathbb{N}}
\newcommand*\bP{\mathbb{P}}
\newcommand*\bQ{\mathbb{Q}}
\newcommand*\bR{\mathbb{R}}
\newcommand*\bS{\mathbb{S}}
\newcommand*\bZ{\mathbb{Z}}
\newcommand*\bB{\mathbb{B}}
\newcommand*\bOne{{\mathchoice {\rm 1\mskip-4mu l} {\rm 1\mskip-4mu l} {\rm 1\mskip-4.5mu l}{\rm 1\mskip-5mu l}}}
\newcommand\bDel{\mbox{$\Delta$\newline \scalebox{.6}{\raisebox{3pt}{$\mkern-22mu\not{}\mkern12mu$}}}}
\newcommand\bDelta{\mathchoice{\bDel}{\bDel}{\scalebox{.7}{\bDel}}{\scalebox{.5}{\bDel}}}
\let\leq\leqslant
\let\geq\geqslant
\newcommand*\bydef{\coloneqq}
\newcommand*{\mk}{\mkern -1mu}
\newcommand*{\Mk}{\mkern -2mu}
\newcommand*{\mK}{\mkern 1mu}
\newcommand*{\MK}{\mkern 2mu}
\newcommand*{\romanenumi}{\renewcommand*{\theenumi}{\roman{enumi}}}
\newcommand*{\Romanenumi}{\renewcommand*{\theenumi}{\Roman{enumi}}}
\newcommand*{\alphenumi}{\renewcommand*{\theenumi}{\alph{enumi}}}
\newcommand*{\Alphenumi}{\renewcommand*{\theenumi}{\Alph{enumi}}}
\let\oldtilde\tilde
\renewcommand*{\tilde}[1]{\mathchoice{\widetilde{#1}}{\widetilde{#1}}{\oldtilde{#1}}{\oldtilde{#1}}}
\let\oldexists\exists
\renewcommand*{\exists}{\mathrel{\oldexists}}
\let\oldforall\forall
\renewcommand*{\forall}{\mathrel{\oldforall}}
\newtheorem{theo}{Theorem}[section]
\newtheorem{lemm}[theo]{Lemma}
\newtheorem{prop}[theo]{Proposition}
\newtheorem{coro}[theo]{Corollary}
\theoremstyle{definition}\newtheorem{defi}[theo]{Definition}
\theoremstyle{remark}\newtheorem{rema}[theo]{Remark}
\newtheorem{exam}[theo]{Example}

\title[Chern form inequalities]{On the existence of logarithmic and orbifold jet differentials: original Lemma 6.5}
\author{Fr\'ed\'eric Campana, Lionel Darondeau, Jean-Pierre Demailly and Erwan Rousseau}
\date{Source: Annales Henri Lebesgue 7 (2024), publisher version; DOI 10.5802/ahl.197}
\begin{document}\maketitle
\noindent\textit{Editorial scope.} The counted unit is the original published Lemma 6.5, the sharp two-sided spherical product-moment bound for $p$ non-zero complex linear forms on $\bC^r$ together with its proportionality and orthogonality equality cases, reproduced exactly from native source0.tex lines 1594-1616, followed by its complete author proof reproduced exactly from native lines 1618-1755. Retained as uncounted necessary context: the printed section 6 heading (native line 1399) and the statement of the later Proposition 6.7 (native lines 1771-1805), which is the definition site of the enumerated item that the counted proof cross-refers to at native lines 1694, 1700 and 1707; that proposition is not counted and its own proof is not retained. The subsection heading 6.2 is editorial, because the native line 1593 carries the connector sentence and its forward cross-reference to Section 7 in the same line. Sections 1-5 and 7-8, the introduction, Remark 6.6, the proof of Proposition 6.7 and the remainder of the paper are not selected. Original printed slips are preserved literally without repair: the missing verb in the statement phrase "and let $\mu$ the rotation invariant probability measure"; the polar-coordinate line at native 1624 which writes $u^\prime\in\bS^{2r-1}$ where the disintegration uses $\bS^{2r-3}$; and the three cross-references at native 1694, 1700 and 1707 which resolve to the enumerated item (b) of the later Proposition 6.7 rather than to item (b) of Lemma 6.5 itself, so that the published print "6.5(b)" is correct only by coincidence of both items being labelled (b).
\setcounter{section}{5}
\section{Positivity concepts for vector bundles and Chern inequalities}\label{se:Chern_inequalities}

\setcounter{subsection}{1}
\subsection{Chern form inequalities}
\setcounter{theo}{4}
\begin{lemm}\label{lemm:Chern_form_inequalities}
Let $\ell_j\in(\bC^r)^*$, $1\leq j\leq p$, be non-zero complex linear forms on $\bC^r$, where $(\bC^r)^*\simeq\bC^r$ is equipped with its standard hermitian form, and let $\mu$ the rotation invariant probability measure on $\bS^{2r-1}\subset\bC^r$. Then
\[
I(\ell_1,\,\dotsc,\,\ell_p)=\int_{\bS^{2r-1}}\abs{\ell_1(u)}^2\ldots\,\abs{\ell_p(u)}^2\, d\mu(u)
\]
satisfies the following inequalities:\\
\begin{eqlist}
\[
\tag{a}\label{item:Chern_form_inequality1}
I(\ell_1,\,\dotsc,\,\ell_p)
\le
\tfrac{p!\,(r-1)!}{ (p+r-1)!}\,\prod_{j=1}^p\abs{\ell_j}^2,
\]
and the equality occurs if and only if the $\ell_j$ are proportional;
\[
\tag{b}\label{item:Chern_form_inequality2}
I(\ell_1,\,\dotsc,\,\ell_p)
\ge
\tfrac{(r-1)!}{ (p+r-1)!}\,\prod_{j=1}^p\abs{\ell_j}^2,
\]
and the equality occurs if and only if $p\leq r$ and the $\ell_j$ are pairwise orthogonal.
\end{eqlist}
\end{lemm}

\begin{proof}
Denote by $d\lambda$ the Lebesgue measure on Euclidean space and by $d\sigma$ the area measure of the sphere. One can easily check that the projection
\[
\bS^{2r-1}\to \bB^{2r-2},\quad u=\left(u_1,\,\dotsc,\,u_r\right)\mapsto v=\left(u_1,\,\dotsc,\,u_{r-1}\right),
\]
yields $d\sigma(u)=d\theta\wedge d\lambda(v)$ where $u_r=\abs{u_r}\,e^{i\theta}$ [just check that the wedge products of both sides with $\tfrac{1}{ 2}d\abs{u}^2$ are equal to $d\lambda(u)$, and use the fact that $d\theta=\tfrac{1}{ 2i}(du_r/u_r -d\bar u_r/\bar u_r)$], thus, in terms of polar coordinates $v=t\,u'$, $u'\in\bS^{2r-1}$, we have
$d\sigma(u)=d\theta\wedge t^{2r-3}\,dt\wedge d\sigma'(u')$, and going back to the invariant probability measures $\mu$ on $\bS^{2r-1}$ and $\mu'$ on $\bS^{2r-3}$, we get
$\abs{u_r}^2=1-\abs{v}^2=1-t^2$ and an equality
\begin{equation}\label{eq:measure}
d\mu(u)=\tfrac{2r-2}{ 2\pi}\,d\theta\wedge t^{2r-3}\,dt
\wedge d\mu'(u').
\end{equation}
If $\ell_1,\,\dotsc,\,\ell_p$ are independent of $u_r$, \eqref{eq:measure} and the Fubini theorem imply by homogeneity
\begin{align}\label{5.10}
\int_{\bS^{2r-1}}\abs{\ell_1(u')}^2\ldots\,\abs{\ell_p(u')}^2\,d\mu(u) =
\tfrac{r-1}{p+r-1}
\int_{\bS^{2r-3}}\abs{\ell_1(u')}^2\ldots\,\abs{\ell_p(u')}^2\,d\mu'(u'),
\end{align}
\begin{multline}\label{5.10'}
\tag{\theequation\ensuremath{'}}
\int_{\bS^{2r-1}}\abs{\ell_1(u')}^2\ldots\,\abs{\ell_{p-1}(u')}^2\,\abs{u_r}^2\,d\mu(u) =
\\
\tfrac{r-1}{(p+r-2)(p+r-1)} \int_{\bS^{2r-3}}\abs{\ell_1(u')}^2\ldots\,\abs{\ell_{p-1}(u')}^2\,d\mu'(u')
\end{multline}
(for instance, in case~\eqref{5.10'}, we have to integrate $t^{2p-2}(1-t^2)\times t^{2r-3}\,dt$). For $p\leq r$, the formulas
\[
\int_{\bS^{2r-1}}\abs{u_1}^{2p}\,d\mu(u)=\tfrac{p!\,(r-1)!}{ (p+r-1)!},\quad
\int_{\bS^{2r-1}}\abs{u_1}^2\ldots\,\abs{u_p}^2\,d\mu(u)=\tfrac{(r-1)!}{ (p+r-1)!},
\]
are then obtained by induction on $r$ and $p$.
\begin{enumerate}
\item For any $\ell\in(\bC^r)^*$, we can find orthonormal coordinates on $\bC^r$ such that $\ell(u)=\abs{\ell}\,u_1$ in the new coordinates. Hence
\[
\int_{\bS^{2r-1}}\abs{\ell(u)}^{2p}\,d\mu(u)=m_{r,p}\,\abs{\ell}^{2p}\quad
\text{where}~~m_{r,p}=\int_{\bS^{2r-1}}\abs{u_1}^{2p}\,d\mu(u)=
\tfrac{p!\,(r-1)!}{ (p+r-1)!}.
\]
It follows from H\"older's inequality that
\[
I(\ell_1,\,\dotsc,\,\ell_p)\leq \prod_{j=1}^p
\left(\int_{\bS^{2r-1}}\abs{\ell_j}^{2p}\,d\mu(u)\right)^{1/p}=m_{r,p}
\prod_{j=1}^p\abs{\ell_j}^2,
\]
and that the equality occurs if and only if all $\ell_j$ are proportional.

\item We prove the inequality
\[
I(\ell_1,\,\dotsc,\,\ell_p)\geq \tfrac{(r-1)!}{(p+r-1)!}\prod_{j=1}^p\abs{\ell_j}^2
\]
by induction on $p$, the result being clear for $p=0$ or $p=1$. If we choose an orthonormal basis $(e_1,\,\dotsc,\,e_r)\in\bC^r$ such that
$\ell_j(e_r)\ne 0$ for all $j$ and replace $\ell_j$ by
$(\ell_j(e_r))^{-1}\ell_j$, we can assume $\ell_j(e_r)=1$. We then write $u=u'+u_re_r$ with $u'\in e_r^\perp\simeq\bC^{r-1}$ and
\[
\ell_j(u)=\ell_j'(u')+u_r,\quad 1\leq j\leq p,\quad\ell'_j
\bydef
\ell_{j|e_r^\perp}.
\]
Let $s_k(\ell'_\bullet(u'))$ be the elementary symmetric functions in $\ell'_j(u')$, $1\leq j\leq p$, with $s_0\bydef1$. We have
\[
I(\ell_1,\,\dotsc,\,\ell_p) =\int_{\bS^{2r-1}}\prod_{j=1}^p\left|\ell'_j(u')+u_r\right|^2\,d\mu(u) =\int_{\bS^{2r-1}}
\abs*{\sum_{k=0}^ps_k(\ell'_\bullet(u'))\,u_r^{p-k}}^2 d\mu(u).
\]
We make a change of variable $u_r\mapsto u_r\,e^{i\theta}$ and take the average over $\theta\in[0,2\pi]$. Parseval's formula gives
\[
I(\ell_1,\,\dotsc,\,\ell_p) =\int_{\bS^{2r-1}}\sum_{k=0}^p\abs*{s_k(\ell'_\bullet(u'))}^2\,\abs{u_r}^{2(p-k)} d\mu(u),
\]
and since
\[
(2r-2)\int_0^1t^{2k}(1-t^2)^{p-k}\,t^{2r-3}dt=
\tfrac{(r-1)\,(k+r-2)!\,(p-k)!}{(p+r-1)!},
\]
formula~\eqref{eq:measure} implies
\[
I(\ell_1,\,\dotsc,\,\ell_p) =\int_{\bS^{2r-3}}\sum_{k=0}^p\tfrac{(r-1)\,(k+r-2)!\,(p-k)!}{(p+r-1)!}\,
\abs*{s_k(\ell'_\bullet(u'))}^2\,d\mu'(u').
\]
As $\abs{\ell_j}^2=1+\abs{\ell'_j}^2$, our inequality~\ref{lemm:Chern_form_inequalities}\eqref{item:Chern_curv_inequality2} is equivalent to
\begin{equation}\label{eq:basic_inequality}
\int_{\bS^{2r-3}}\sum_{k=0}^p\tfrac{(k+r-2)!\,(p-k)!}{(r-2)!}\,
\abs*{s_k(\ell'_\bullet(u'))}^2\,d\mu'(u')\ge
\prod_{j=1}^p\left(1+\abs{\ell'_j}^2\right)
\end{equation}
for all linear forms $\ell'_j\in(\bC^{r-1})^*$. We actually prove~\eqref{eq:basic_inequality} by induction on $p$ (observing that the inequality is a trivial equality for $p=0,1$). Assume that~\eqref{eq:basic_inequality} (and hence~\ref{lemm:Chern_form_inequalities}\eqref{item:Chern_curv_inequality2}) is known for any $(p-1)$-tuple of linear forms $(\ell'_1,\,\dotsc,\,\ell'_{p-1})$. As~\ref{lemm:Chern_form_inequalities}\eqref{item:Chern_curv_inequality2} is invariant under the action of $U(r)$, it is sufficient to consider the case when $\ell_p(u)=u_r$, i.e.\ $\ell'_p=0$. The induction hypothesis tells us that
\[
\int_{\bS^{2r-3}}\sum_{k=0}^{p-1}\tfrac{(k+r-2)!\,(p-1-k)!}{(r-2)!}\,
\abs*{s_k(\ell'_\bullet(u'))}^2\,d\mu'(u')\ge
\prod_{j=1}^{p-1}(1+\abs{\ell'_j}^2).
\]
However, when we add the factor $\ell_p$, the elementary symmetric functions $s_k(\ell'_\bullet(u'))$ are left unchanged for $k\leq p-1$, while
$s_p(\ell'_\bullet(u'))=0$ and $1+\abs{\ell'_p}^2=1$. Therefore~\eqref{eq:basic_inequality} holds true for $p$, since $(p-k)!\geq (p-1-k)!$ for all $k=0,1,\dotsc,p-1$. We have proved the inequality at order $p$ whenever
$\ell_p=\alpha_p\langle\bullet,e_r\rangle$ and $\ell_j(e_r)\ne 0$ for $j\leq p-1$. Since those $(\ell_1,\,\dotsc,\,\ell_p)$ are dense in the space $((\bC^r)^*)^p$ of $p$-tuples of linear forms, the proof of the lower bound is complete.

\item (b, equality case) We argue by induction on $r$. For $r=1$, we have in fact $\ell_j(u)=\alpha_ju_1$, $\alpha_j\in\bC^*$, and
$I(\ell_1,\,\dotsc,\,\ell_r)=\prod\abs{\ell_j}^2$, thus the coefficient $\tfrac{1}{ (p+r-1)!}=\tfrac{1}{ p!}$ is reached if and only if $p\leq 1$. Now, assume $r\geq 2$ and the equality case solved for dimension $r-1$. By rescaling and reordering the $\ell_j$, we can always assume that $\ell_j(e_r)\ne 0$ (and hence
$\ell_j(e_r)=1$) for $q+1\leq j\leq p$, while $\ell_j(e_r)=0$ for
$1\leq j\leq q$ (we can possibly have $q=0$ here). Then we write
$\ell_j(u)=\ell'_j(u')$ for $1\leq j\leq q$ and
$\ell_j(u)=\ell'_j(u')+u_r$ for $q+1\leq j\leq p$. Therefore, if
$s_k(\ell'(u'))$ denotes the $k^{\rm th}$ elementary symmetric function in
$(\ell'_j(u')_{q+1\leq j\leq p}$, we find
\begin{align*}
I(\ell_1,\,\dotsc,\,\ell_p) &=\int_{\bS^{2r-1}}\prod_{j=1}^q\left|\ell'_j(u')\right|^2\,
\prod_{j=q+1}^p\left|\ell'_j(u')+u_r\right|^2\,d\mu(u)\\
&=\int_{\bS^{2r-1}}\prod_{j=1}^q\left|\ell'_j(u')\right|^2
\left|\sum_{k=0}^{p-q}s_k(\ell'(u'))\,u_r^{p-q-k}\right|^2 d\mu(u)
\\
&=\int_{\bS^{2r-1}}\prod_{j=1}^q\left|\ell'_j(u')\right|^2
\sum_{k=0}^{p-q}
\left|s_k(\ell'(u'))\right|^2\,\abs{u_r}^{2(p-q-k)}\,d\mu(u)\\
&=\int_{\bS^{2r-3}}\prod_{j=1}^q\left|\ell'_j(u')\right|^2
\sum_{k=0}^{p-q}\tfrac{(r-1)\,(k+r-2)!\,(p-q-k)!}{(p-q+r-1)!}\,
\left|s_k(\ell'(u'))\right|^2\,d\mu'(u')\\
&\geq\tfrac{(r-1)!}{(p+r-1)!}\,\prod_{j=1}^q\abs{\ell'_j}^2
\prod_{j=q+1}^p\left(1+\abs{\ell'_j}^2\right)
\end{align*}
by what we have just proved. In an equivalent way, we get
\begin{multline*}
\int_{\bS^{2r-3}}\prod_{j=1}^q\left|\ell'_j(u')\right|^2
\sum_{k=0}^{p-q}\tfrac{(k+r-2)!\,(p-q-k)!\,(p+r-1)!}{(r-2)!\,(p-q+r-1)!}\,
\left|s_k(\ell'(u'))\right|^2\,d\mu'(u')
\\
\geq\prod_{j=1}^q\abs{\ell'_j}^2
\prod_{j=q+1}^p\left(1+\abs{\ell'_j}^2\right)
\end{multline*}
In general, we can rotate coordinates in such a way that $\ell_p(u)=u_r$ and $\ell'_p=0$, and we see that the above inequality holds when $p$ is replaced by $p-1$, as soon as $q\leq p-2$. Then the corresponding coefficients $k=0$ for $p$, $p-1$ are
\[
\tfrac{(p-q)!\,(p+r-1)!}{(p-q+r-1)!}>\tfrac{(p-1-q)!\,(p-1+r-1)!}{(p-1-q+r-1)!},
\]
and since $s_0=1$, we infer that the inequality is strict. The only possibility for the equality case is $q=p-1$, but then
\[
I(\ell_1,\,\dotsc,\,\ell_p)=
\int_{\bS^{2r-1}}\prod_{j=1}^{p-1}\left|\ell'_j(u')\right|^2\,\abs{u_r}^2\,d\mu(u)=
\tfrac{r-1}{ p+r-1}\int_{\bS^{2r-3}}\prod_{j=1}^{p-1}\left|\ell'_j(u')\right|^2\,d\mu'(u'),
\]
and we see that we must have equality in the case $(r-1,p-1)$. By induction, we conclude that $p-1\leq r-1$ and that the
$\ell_j(u)=\ell'_j(u')$ are orthogonal for $j\leq p-1$, as desired.\qedhere
\end{enumerate}
\end{proof}

\setcounter{theo}{6}
\begin{prop}\label{prop:Chern_curv_inequalities}
Let $T$, $E$ be complex vector spaces of respective dimensions $\dim T=n$, $\dim E=r$. Assume that $E$ is equipped with a hermitian structure $h$, and denote by $\mu$ the unitary invariant probability measure $\mu$ on the unit sphere bundle $S(E)=\{u\in E\,;\abs{u}_h=1\}$ of $E$.
\begin{enumerate}\alphenumi
\item\label{item:Chern_curv_inequality1}
If $\ell_1,\,\dotsc,\,\ell_k\in E^*$ and $\theta_1,\dotsc,\theta_{p-k}\geq_S 0$ are strongly semi-positive hermitian tensors in
$\Herm(T\otimes E)\simeq\Lambda_\bR^{1,1}T^*\otimes_\bR\Herm(E,E)$, then
\begin{multline*}
\int_{u\,\in\,S(E)}\abs{\ell_1(u)}^2\ldots\abs{\ell_k(u)}^2\,
\langle\theta_1(u),u\rangle_h\wedge\ldots
\wedge\langle\theta_{p-k}(u),u\rangle_h\,d\mu(u)
\\
\begin{cases}
\displaystyle
\geq \tfrac{(r-1)!}{ (p+r-1)!}\,\bigg(\prod_{j=1}^k\abs{\ell_j}^2\bigg)
\Tr_h\theta_1\wedge\ldots\wedge\Tr_h\theta_{p-k},\\
\displaystyle
\leq \tfrac{p!\,(r-1)!}{ (p+r-1)!}\,\bigg(\prod_{j=1}^k\abs{\ell_j}^2\bigg)
\Tr_h\theta_1\wedge\ldots\wedge\Tr_h\theta_{p-k},
\end{cases}
\end{multline*}
as pointwise strong inequalities of $(p-k,p-k)$-forms.
\item\label{item:Chern_curv_inequality2}
If $\theta\geq_G 0$ in $\Lambda_\bR^{1,1}T^*\otimes_\bR\Herm(E,E)$ and $\ell_j\in E^*$, then
\[
\int_{u\,\in\,S(E)}\abs{\ell_1(u)}^2\ldots\abs{\ell_k(u)}^2\,
\langle\theta(u),u\rangle_h^{p-k}\,d\mu(u)\le
\tfrac{p!\,(r-1)!}{ (p+r-1)!}\bigg(\prod_{j=1}^k\abs{\ell_j}^2\bigg)(\Tr_h\theta)^{p-k}
\]
as a pointwise weak inequality of $(p-k,p-k)$-forms.
\end{enumerate}
In particular, the above inequalities apply when
$(E,h)$ is a hermitian holomorphic vector bundle of rank $r$ on a complex $n$-dimensional manifold $X$, and one takes
$\theta_j=\Theta_{E,h}$ to be the curvature tensor of $E$, so that
$\Tr_h\theta_j=c_1(E,h)$ is the first Chern form of $(E,h)$.
\end{prop}

\end{document}
